\section{Conclusion and Future Work}

We built a controlled two-track experiment for measuring transfer under four
nested label budgets: one detector, one corrected label contract, one
strict-FP32 recipe, and a fail-closed path from training bytes to every
published number. The completed 32-cell grid supports three conclusions.
First, transfer concentrates its value under scarcity---at 12 scenes every
pretrained arm except CNN-optical clears its random floor, and the best
12-scene cells approach what 111 scenes buy a randomly initialized
detector. Second, matching the sensing modality is not uniformly the right
scarce-label choice: SAR pretraining dominates optical in the CNN track at
every budget, while the ViT track prefers the optical checkpoint below half
data, so the contrast's sign is architecture-dependent. Third, generic
ImageNet initialization wins both tracks at full data. These remain
corrected development-selection results from one seed and support no
variance claim.

\ifHevTestComplete
The once-scored 16-scene test column bounds development-selection optimism
at a mean change of \HevTestMeanGap{} F1 without altering the within-track
ordering conclusions.
\else
The next step scores the frozen 16-scene test split under the held-out
contract, which applies each cell's development-bound threshold verbatim.
\fi
\ifHevFinalEval
The once-only 50-scene human-verified evaluation completes the dark-vessel
and near-shore analysis.
\else
The 50 human-verified scenes then support one sealed evaluation of
dark-vessel recall and near-shore performance for follow-on work.
\fi

\paragraph{Roles.} John Roth and Kyle Wagner jointly designed the study,
implemented its data, model, training, and evaluation code, ran the
experiments, analyzed the evidence, and wrote this report.
